% SPDX-License-Identifier: LPPL-1.3c
% Copyright (C) 2026 Christophe Jorssen
% Author and Current Maintainer: Christophe Jorssen <christophe.jorssen@gmail.com>
% LPPL maintenance status: maintained
% This file is part of luafemm. See LICENSE and LICENSES.md for its terms.
\section{Saving and reusing computations}
\label{sec:cache}
The mesh and converged solution can be reused across compilations. This is
particularly useful while editing labels, field-line styles or a PGFPlots
axis. The document still declares its physical paths: they identify the
requested calculation and let luafemm detect outdated results.

\begin{codeexample}[width=5cm]
\begin{tikzpicture}[
  femm/problem={xmin=-20,xmax=20,ymin=-20,ymax=20,mesh size=2},
  femm/cache={file=magnet-field,mode=auto}]
  \filldraw[femm/region={material=cast-alnico-5-lng37,
    magnetization angle=90},fill=red!20]
    (-4,-6) rectangle (4,6);
  \draw[femm/profile={name=probe,samples=101},dashed]
    (-12,9)--(12,9);
  \pic[femm/lines=17] {femm field};
\end{tikzpicture}
\end{codeexample}

On the first compilation the field request meshes and solves the problem,
then saves the result. Later compilations restore the nodal potential and
derive the element fields from the current material declarations. Profiles
and field lines are regenerated on demand, so their graphical and sampling
options remain independent of the saved calculation.

\subsection{Picture options}
\begin{key}{/tikz/femm/cache=\marg{options} (default \normalfont empty)}
Configure persistent reuse for a model. Like |femm/problem|, this key
belongs in the picture options, before the model's begin-picture hook
runs. It may precede or follow |femm/problem|. Without this key, caching
is off. Invoking it sets |mode=auto| before processing its subkeys.
Thus |femm/cache| alone enables automatic reuse with the default filename.
Normal PGF grouping and inherited styles apply.
\end{key}
\begin{key}{/femm/cache/file=\meta{filename stem} (initially \string\jobname-femm)}
The package appends |.lfc|. For example, |file=results/u-magnet| uses
\nolinkurl{results/u-magnet.lfc}. Macros are expanded and passed as escaped
strings. Spaces in filenames are supported.

Relative names are resolved under TeX's output directory when one is set,
otherwise under the compilation's working directory. With
|--output-directory=build|, |file=magnet| therefore writes
\nolinkurl{build/magnet.lfc}. Absolute names are used as supplied.
Any parent directory must already exist; luafemm does not create directories.
Use a distinct stem for each physical problem and for concurrent writers.
The default is convenient for a document containing one magnetic model.
\end{key}
\begin{key}{/femm/cache/mode=\meta{choice} (initially off)}
A PGF |.is choice| selector. Unknown choices are errors. The enclosing
|femm/cache| key selects |auto| unless explicitly overridden.
\end{key}
\begin{key}{/femm/cache/mode/auto}
Restore a matching solution. If only the mesh matches, restore and
reclassify it with the current materials and excitations, then solve.
Otherwise generate a new mesh and solve. Missing, damaged or incompatible
records produce a logged miss and are replaced after successful preparation.
\end{key}
\begin{key}{/femm/cache/mode/refresh}
Ignore disk contents, generate a new mesh and solution when requested,
and replace the cache. Repeated requests within the same model reuse its
in-memory data; refresh does not mean solving for every field picture.
\end{key}
\begin{key}{/femm/cache/mode/off}
Neither read nor write cache files. Existing files are left untouched.
This is the behaviour of documents that do not enable |femm/cache|.
\end{key}
\begin{key}{/femm/cache/mode/frozen}
Require compatible saved data and never write a file. A mesh picture
requires a matching mesh; a field picture, explicit solve or sampled
profile requires a matching converged solution. A missing, damaged,
incompatible or outdated record raises an error before numerical work.
It never silently displays an outdated field or falls back to a solve.
Changing presentation options remains allowed.
\end{key}

\subsection{What invalidates the result?}
\begin{center}
\begin{tabular}{p{.57\linewidth}p{.34\linewidth}}
\toprule Change & Work in |auto| mode \\
\midrule
Domain, contours, fill rule, units, mesher, spacing, mesh tolerance or quality controls
  & Generate mesh and solve \\
Material law, coercivity, region current, magnetisation or boundary data
  & Reuse mesh; solve again \\
Solver tolerance or iteration limit & Reuse mesh; solve again \\
Colour, stroke, arrows, field-line count, clip, labels
  & Reuse solution \\
Profile path, sample count, component, offset or PGFPlots options
  & Reuse solution; sample again \\
Incompatible package or cache-format version & Generate mesh and solve \\
\bottomrule
\end{tabular}
\end{center}

The geometry descriptor uses the evaluated, flattened and rounded model
coordinates, not the source text. A curve-tolerance change invalidates the
mesh when it changes those coordinates. Whole-picture transformations
normally affect only presentation; transformations inside a physical path
affect the model. The current picture frame is always retained. Precision
effects that change captured coordinates can cause a conservative cache miss.
Region order matters because later filled regions take material precedence.
Every subcontour and the native fill rule are included, even when a hole
contains only background air. A fill-rule change can alter local refinement
and therefore invalidates the mesh as well as the solution.

The solution descriptor contains the actual constitutive data, including
B--H points and effective coercivity, rather than just catalogue identifiers.
Unused materials, log callbacks, frames and profile definitions are excluded.
Descriptors are compared in full; a short checksum collision cannot cause a
stale model to match. This is a single-record cache, not a history of parameter
sets: changing a model back may require another solve.

\subsection{Diagnostics and saved statistics}
\begin{command}{\femmcachestatus}
Print the current model's cache outcome without starting a calculation:
|off|, |pending|, |miss|, |mesh|, |solution| or |refresh|.
|pending| means caching is enabled but no mesh or solve has been requested.
|mesh| means only the mesh was restored; after an automatic solve it still
records that reuse outcome. |solution| means a converged solution was restored.
|miss| and |refresh| describe how the current calculation was obtained, even
after it has been saved successfully. A current model must exist.
\end{command}

The log gives the resolved filename and the reason for a miss or partial
reuse. Lua scripts can inspect |model.cache_info|. A restored solution
retains the original Newton/CG counts, residual and solution time in
|\femmstat|; the mesh diagnostics likewise describe its original generation.
In particular, |seconds| is \emph{not} the current cache-loading time.
All material and region declarations must precede the first mesh or solve.

\subsection{File handling and scope}
The private |.lfc| format stores passive, versioned data with lossless
binary64 round trips: geometry and physical descriptors, SI mesh vertices,
connectivity, exterior tags, diagnostics and, when available, nodal potential.
It is parsed without |load|, |loadfile| or |dofile|. Size, integrity and
structural checks run before restoration. A checksum detects accidental
damage; it is not an authentication mechanism.

Writes use a neighbouring temporary file and a final rename. A failed write
or replacement leaves the previous target intact and raises an error. If
the filesystem cannot replace an existing target through rename, the
operation fails explicitly. File-access errors on writing are not silently
ignored. Only a converged solution is saved; an unsuccessful solve may
leave a valid mesh-only cache. Interrupted processes can leave temporary
files, which are not considered cache records.

The current format supports both luafemm meshers and can be shared among
plain LuaTeX, LuaLaTeX and ConTeXt MkIV using the same package version.
Lua-only source and boundary closures require caching to be off, because
their captured state cannot be fingerprinted reliably.

These files are disposable computational caches, not FEMM |.fem| or |.ans|
files and not a replacement for the document source. FEMM exchange and
loading a foreign model without TeX declarations are outside this version.
Keep the source, material data and package version to reproduce a result;
use |frozen| to prevent an unintended recomputation while editing.
